\section{Relative generation by continuity}
\label{sec:relative-generation}

We now convert the correlation supplied by Proposition~\ref{prop:perturbation}
into generation. Throughout this section, $P\subset M$ is an irreducible
inclusion of $\mathrm{II}_1$ factors and $M$ has separable predual.

\begin{lemma}
\label{lem:unitary-completeness}
The Hilbert space $L^2(M,\tau)$ is separable, and
$\mathcal U(M)$ is complete for the metric $d(u,v)=\norm{u-v}_2$.
\end{lemma}

\begin{proof}
The unit ball of $M$ is compact and metrizable in the weak-star topology,
so it has a countable weak-star dense subset $S$. For each
$\xi\in L^2(M,\tau)$, the functional
$a\mapsto\langle a,\xi\rangle$ on $M$ is normal. Indeed, normal
functionals form a norm-closed subspace of $M^*$, and approximation
of $\xi$ by elements of $M$ in $L^2$ gives norm approximation of
this functional by normal functionals, since
\[
 \sup_{\norm{a}\leq1}
 |\langle a,\xi-\eta\rangle|
 \leq\norm{\xi-\eta}_2.
\]
If $\xi$ is orthogonal to $S$, normality implies that it is orthogonal to
the unit ball, and hence to $M$. Thus $S$ has dense linear span in
$L^2(M,\tau)$, proving separability.

Let $(u_n)$ be a Cauchy sequence in $\mathcal U(M)$ for $d$.
The adjoints are also Cauchy, because
$\norm{u_n^*-u_m^*}_2=\norm{u_n-u_m}_2$.
For $a\in M$,
\[
 \norm{(u_n-u_m)a}_2\leq\norm{a}\,\norm{u_n-u_m}_2.
\]
The left multiplication operators associated with $u_n$ therefore
converge strongly on the dense subspace $M\subset L^2(M,\tau)$.
Their uniform operator norm bound extends this convergence to the whole
Hilbert space. The same holds for the left multipliers associated with
$u_n^*$. Strong closedness of the left regular representation of $M$
gives limits $u,v\in M$, and passage to inner products gives $v=u^*$.
Multiplication is jointly strongly continuous on bounded sets, so
$uv=vu=1$. Finally, convergence on the vector $1$ gives
$\norm{u_n-u}_2\to0$.
\end{proof}

For $u\in\mathcal U(M)$, write
\[
 N_u=W^*(P,u),\qquad
 e_u:L^2(M,\tau)\longrightarrow L^2(N_u,\tau)
\]
for the orthogonal projection. For $\xi\in L^2(M,\tau)$, put
$f_\xi(u)=\norm{e_u\xi}_2$.

If $N_u$ is proper, there is a nonzero bounded vector orthogonal to
$L^2(N_u,\tau)$. The perturbation proposition will keep its projection
norm bounded away from zero along a sequence of unitaries converging to
$u$, so we seek points at which every $f_\xi$ is continuous simultaneously.

\begin{lemma}
\label{lem:projection-continuity}
Each $f_\xi$ is bounded and lower semicontinuous on
$(\mathcal U(M),d)$. There is a dense $G_\delta$ subset
$\mathcal C\subset\mathcal U(M)$ such that every $f_\xi$,
$\xi\in L^2(M,\tau)$, is continuous at every point of $\mathcal C$.
\end{lemma}

\begin{proof}
Let $\mathcal P_P$ denote the family of finite $*$-polynomials in a
unitary variable with coefficients in $P$. For each fixed $u$, the
values $b(u)$, $b\in\mathcal P_P$, are dense in $L^2(N_u,\tau)$.
To see this, take the operator norm closure of the polynomial algebra.
Its von Neumann closure is $N_u$, and Kaplansky density approximates
each element of $N_u$ strongly by bounded elements of this
$C^*$-algebra. Evaluation on $1\in L^2(M,\tau)$ gives $2$-norm
approximation, which can then be followed by operator norm
approximation by polynomials. The density of $N_u$ in $L^2(N_u,\tau)$
finishes the assertion.

For fixed $b\in\mathcal P_P$, the map $u\mapsto b(u)$ is continuous
in $2$-norm. This follows by telescoping products, using the fixed
operator norm bounds of the coefficients, the fact that unitaries
have norm one, and the identity
$\norm{u^*-v^*}_2=\norm{u-v}_2$. Consequently
\[
 d_\xi(u):=\operatorname{dist}\bigl(\xi,L^2(N_u,\tau)\bigr)
 =\inf_{b\in\mathcal P_P}\norm{\xi-b(u)}_2
\]
is upper semicontinuous: an infimum of any family of continuous
real functions is upper semicontinuous. Orthogonality gives
\[
 f_\xi(u)^2=\norm{\xi}_2^2-d_\xi(u)^2,
\]
so $f_\xi$ is lower semicontinuous and
$0\leq f_\xi(u)\leq\norm{\xi}_2$.

For completeness, a bounded lower semicontinuous real function $f$
on a complete metric space is continuous on a dense $G_\delta$ set.
For $\varepsilon>0$, let $G_\varepsilon$ be the union of all open
sets on which the oscillation of $f$ is less than $\varepsilon$.
This union is dense. Indeed, in any nonempty open set $O$, write
$s=\sup_O f$ and choose $x\in O$ with
$f(x)>s-\varepsilon/3$. Lower semicontinuity gives an open
neighborhood $V\subset O$ of $x$ on which
$f>f(x)-\varepsilon/3$. The oscillation on $V$ is less than
$\varepsilon$. Baire's theorem now shows that
$\bigcap_{n\geq1}G_{1/n}$ is dense, and $f$ is continuous at each
of its points.

Apply this construction to $f_\eta$ for $\eta$ in a countable dense
subset of $L^2(M,\tau)$ and intersect the resulting dense
$G_\delta$ sets. By Baire's theorem, their intersection
$\mathcal C$ is again a dense $G_\delta$ set.
For arbitrary $\xi,\eta\in L^2(M,\tau)$, contractivity of $e_u$ gives
\begin{equation}
 \label{eq:projection-vector-lipschitz}
 |f_\xi(u)-f_\eta(u)|\leq\norm{\xi-\eta}_2
 \qquad(u\in\mathcal U(M)).
\end{equation}
Thus, for $u\in\mathcal C$ and $v\in\mathcal U(M)$,
\[
 |f_\xi(v)-f_\xi(u)|
 \leq2\norm{\xi-\eta}_2+|f_\eta(v)-f_\eta(u)|.
\]
Approximating $\xi$ by an $\eta$ from the countable dense subset,
and then letting $v\to u$, proves continuity of every $f_\xi$ at $u$.
\end{proof}

\begin{proof}[Proof of Theorem~\ref{thm:relative-generation}]
Take $u\in\mathcal C$ from Lemma~\ref{lem:projection-continuity},
and suppose that $N_u\ne M$. There is a unitary $z\in M\setminus N_u$,
since the unitaries linearly span $M$. Let $E_{N_u}:M\to N_u$ be the
trace-preserving conditional expectation, and put
\[
 \xi=z-E_{N_u}(z)\in M,\qquad s=\norm{\xi}_2>0.
\]
The conditional expectation is the restriction to $M$ of $e_u$.
Therefore
\[
 e_u\xi=0,\qquad f_\xi(u)=0,\qquad
 \tau(\xi^*z)=s^2.
\]
Fix
\begin{equation}
 \label{eq:fixed-alpha-generation}
 0<\alpha<\min\{1/4,s/32\}.
\end{equation}
For every integer $\ell\geq1$, Proposition~\ref{prop:perturbation}
provides a unitary $u_\ell\in M$ and unitaries
$t_{0,\ell},\ldots,t_{\ell,\ell}\in P$ such that, with
\[
 W_\ell=t_{0,\ell}\prod_{k=1}^\ell(u_\ell t_{k,\ell})\in\mathcal U(N_{u_\ell}),
\]
one has
\begin{align}
 \norm{u_\ell-u}_2
 &\leq\frac{\alpha}{\sqrt{2\ell}}+\frac1\ell,
 \label{eq:generation-perturbation-size}\\
 \left|\tau(\xi^*W_\ell)-\frac{\alpha s^2}{2}\right|
 &\leq s\frac{4\alpha^2}{1-2\alpha}+\frac1\ell.
 \label{eq:generation-correlation}
\end{align}
The fixed choice~\eqref{eq:fixed-alpha-generation} gives
\[
 s\frac{4\alpha^2}{1-2\alpha}
 <8s\alpha^2<\frac{\alpha s^2}{4}.
\]
It follows from~\eqref{eq:generation-correlation} that
\[
 |\tau(\xi^*W_\ell)|\geq\frac{\alpha s^2}{4}-\frac1\ell
 \geq\frac{\alpha s^2}{8}
\]
for all sufficiently large $\ell$. On the other hand,
$\norm{W_\ell}_2=1$ and $W_\ell\in L^2(N_{u_\ell},\tau)$, so
\[
 |\tau(\xi^*W_\ell)|
 =|\langle W_\ell,e_{u_\ell}\xi\rangle|
 \leq f_\xi(u_\ell).
\]
By~\eqref{eq:generation-perturbation-size}, $u_\ell\to u$ in
$2$-norm. Continuity of $f_\xi$ at $u$ gives
$f_\xi(u_\ell)\to f_\xi(u)=0$, a contradiction. Thus
$M=W^*(P,u)$ for every $u\in\mathcal C$.

To see that the full set of relative generators is $G_\delta$, choose
$(\eta_j)_{j\geq1}$ dense in $L^2(M,\tau)$. With $d_{\eta_j}$ as in
Lemma~\ref{lem:projection-continuity}, this set is
\[
 \bigcap_{j,n\geq1}\{v\in\mathcal U(M):d_{\eta_j}(v)<1/n\}.
\]
Each set in the intersection is open by upper semicontinuity.
Vanishing of every distance is equivalent to
$L^2(N_v,\tau)=L^2(M,\tau)$, hence to $N_v=M$ by the
trace-preserving conditional expectation. The set is dense because
it contains $\mathcal C$.
\end{proof}

\section{Two self-adjoint generators}
\label{sec:two-generators}

Relative generation over an irreducible hyperfinite subfactor will first
give four self-adjoint generators for every separable $\mathrm{II}_1$
factor. We then apply this
uniform bound to a half-trace corner $N$ of $M$: the identification
$M\cong M_2(N)$ allows the four generators of $N$ to be encoded by two
self-adjoint matrices.

We use Popa's irreducible hyperfinite embedding theorem in the
following form.

\begin{theorem}[Popa]
\label{thm:irreducible-hyperfinite}
Let $N\subset M$ be an irreducible inclusion of $\mathrm{II}_1$
factors with separable preduals. There is a hyperfinite
$\mathrm{II}_1$ subfactor $R\subset N$ such that
$R'\cap M=\mathbb C1$.
\end{theorem}

The special case $N=M$ is
\cite[Corollary~4.1]{Popa1981}; it is also used explicitly in
\cite[proof of Theorem~3.8]{Sherman2012}. The relative form above is
recalled in~\cite[Introduction]{Popa2021Ergodic}.

The two-generator result for the hyperfinite factor is classical
\cite{SuzukiSaito1963}. We give a tensor-model proof for the later
count of generators.

\begin{lemma}
\label{lem:hyperfinite-generators}
The hyperfinite $\mathrm{II}_1$ factor with separable predual is
generated by two self-adjoint elements.
\end{lemma}

\begin{proof}
By uniqueness of the hyperfinite $\mathrm{II}_1$ factor with separable
predual \cite[Section~1.3]{Popa2014}, we may use its standard tensor model
\[
 R=\overline{\bigotimes}_{m\geq1}
 \bigl(M_2(\mathbb C),\operatorname{tr}_2\bigr).
\]
On the $m$th tensor site, let
\[
 d_m=\begin{pmatrix}1&0\\0&0\end{pmatrix},\qquad
 q_m=\frac12\begin{pmatrix}1&1\\1&1\end{pmatrix},
\]
with the identity on the other sites. On that site,
$2d_mq_m(1-d_m)$ is the matrix unit $e_{12}$, so $d_m,q_m$
generate the full matrix algebra. Each of the two families
$(d_m)$ and $(q_m)$ consists of commuting projections.

For any countable family of commuting projections $(e_m)$, the
operator norm convergent series
\[
 b=\sum_{m\geq1}2\,3^{-m}e_m
\]
recovers every $e_m$ by spectral calculus. Inductively, after
recovering $e_1,\ldots,e_{m-1}$, form
\[
 r_m=b-\sum_{j<m}2\,3^{-j}e_j.
\]
The restriction of $r_m$ to $1-e_m$ lies between $0$ and $3^{-m}$,
whereas its restriction to $e_m$ lies between $2\,3^{-m}$ and
$3\,3^{-m}$. Since $e_m$ commutes with $r_m$,
\[
 e_m=1_{((3/2)3^{-m},\,\infty)}(r_m)\in W^*(b).
\]
Applying this encoding separately to $(d_m)$ and $(q_m)$ gives
two self-adjoint elements whose generated von Neumann algebra
contains every tensor site, and hence is all of $R$.
\end{proof}

\begin{proposition}
\label{prop:four-generators}
Every $\mathrm{II}_1$ factor with separable predual is generated
by four self-adjoint elements.
\end{proposition}

\begin{proof}
Let $M$ be such a factor. By
Theorem~\ref{thm:irreducible-hyperfinite}, applied with $N=M$, it contains a hyperfinite
$\mathrm{II}_1$ subfactor $P$ with $P'\cap M=\mathbb C1$.
The trace-preserving conditional expectation onto $P$ embeds
$P_*$ isometrically into $M_*$ by precomposition, so $P$ has
separable predual. Lemma~\ref{lem:hyperfinite-generators} gives
self-adjoint $a,b\in P$ with $P=W^*(a,b)$, and
Theorem~\ref{thm:relative-generation} gives $u\in\mathcal U(M)$
with $M=W^*(P,u)$. Hence
\[
 M=W^*\bigl(a,b,\operatorname{Re}u,\operatorname{Im}u\bigr).\qedhere
\]
\end{proof}

The uniform bound in Proposition~\ref{prop:four-generators} also applies
to half-trace corners. The following matrix construction reduces their
four generators to two after amplification. It is the von Neumann algebra
form of the $2\times2$ encoding in~\cite[Proposition~1.1]{DSSW2008},
with the off-diagonal entry fixed to the identity. We include the proof
to make the recovery of the generators and matrix units explicit.

\begin{lemma}
\label{lem:matrix-encoding}
If a unital von Neumann algebra $N$ is generated by four
self-adjoint elements, then $M_2(N)$ is generated by two
self-adjoint elements.
\end{lemma}

\begin{proof}
Write $N=W^*(a_1,a_2,a_3,a_4)$ with $a_j=a_j^*$, and choose
$C>\norm{a_1}+\norm{a_2}$. Put
\[
 A=\begin{pmatrix}a_1&0\\0&a_2+C1_N\end{pmatrix},\qquad
 B=\begin{pmatrix}a_3&1_N\\1_N&a_4\end{pmatrix},\qquad
 Q=W^*(A,B).
\]
Let $E_{ij}$ be the standard matrix units of $M_2(N)$.
The first diagonal block of $A$ has spectrum in
$[-\norm{a_1},\norm{a_1}]$, and the second has spectrum in
$[C-\norm{a_2},C+\norm{a_2}]$. For
$\norm{a_1}<\lambda<C-\norm{a_2}$, spectral calculus therefore gives
\[
 E_{11}=1_{(-\infty,\lambda)}(A)\in Q,
 \qquad E_{22}=1-E_{11}\in Q.
\]
The off-diagonal cuts of $B$ give
$E_{12}=E_{11}BE_{22}\in Q$ and $E_{21}=E_{12}^*\in Q$.
The first corner then contains
\begin{align*}
 a_1E_{11}&=E_{11}AE_{11},&
 a_2E_{11}&=E_{12}AE_{21}-CE_{11},\\
 a_3E_{11}&=E_{11}BE_{11},&
 a_4E_{11}&=E_{12}BE_{21}.
\end{align*}
Thus $NE_{11}\subset Q$. Multiplying this corner by the matrix
units yields $NE_{ij}\subset Q$ for every $i,j$, and hence
$Q=M_2(N)$.
\end{proof}

\begin{proof}[Proof of Theorem~\ref{thm:main}]
Let $M$ be a $\mathrm{II}_1$ factor with separable predual and
normalized trace $\tau$. Choose a projection $p\in M$ with
$\tau(p)=1/2$. Set $N=pMp$, with normalized trace
$\tau_N(a)=2\tau(a)$. This corner is a $\mathrm{II}_1$ factor
with separable predual: precomposition with the normal compression
$a\mapsto pap$ embeds $N_*$ isometrically into $M_*$.
Since $p$ and $1-p$ have the same trace,
there is a partial isometry $v\in M$ with $v^*v=p$ and
$vv^*=1-p$. The corresponding matrix units give a normal
unital $*$-isomorphism
\[
 \Phi:M_2(N)\longrightarrow M,\qquad
 \Phi\begin{pmatrix}b_{11}&b_{12}\\b_{21}&b_{22}\end{pmatrix}
 =b_{11}+b_{12}v^*+vb_{21}+vb_{22}v^*.
\]
Proposition~\ref{prop:four-generators}, applied to $N$, and
Lemma~\ref{lem:matrix-encoding} provide self-adjoint $A,B\in M_2(N)$
with $W^*(A,B)=M_2(N)$. Define
$x=\Phi(A+iB)\in M$. Since the real and imaginary parts of $x$
are $\Phi(A)$ and $\Phi(B)$, respectively,
\[
 W^*(x)=W^*(\Phi(A),\Phi(B))=\Phi(M_2(N))=M.\qedhere
\]
\end{proof}
